\begin{abstract}
% Graph neural networks (GNNs) have attracted extensive research interests for modeling the relational structures of data, ranging from drug design in molecular graphs to anomaly detection in financial networks and beyond. However, in real-world graph data, each node usually possesses multiple labels rather than only one. For instance, in gene interaction networks, each protein may have several functions; in citation networks, papers may belong to multiple domains; and in financial networks, each user might exhibit various types of anomalies. Existing methods can only model each set of labels independently, which on one hand, wastes computational power and time by learning from scratch for each label set, and on the other hand, fails to fully leverage the beneficial correlation implied by the multi-labels, making current methods ineffective in practical multi-label tasks. This paper aims to thoroughly analysis and quantify of the influence between labels across the entire training process. Moreover, based on the analysis and modeling of multi-label relationships, we propose a novel model, \textbf{label information propagation} (\model), specifically designed for multi-label node classification, which represents a strategic advancement by fine-tuning the training dynamics among multiple labels to improve overall performance. Our framework is evaluated through several benchmark datasets, which demonstrates the superior performance beyond state-of-the-art methods.\footnote{Our code is available at \url{https://anonymous.4open.science/r/MultiLblNode} for review.}
Graphs are a complex and versatile data structure used across various domains, with possibly multi-label nodes playing a particularly crucial role. 
Examples include proteins in PPI networks with multiple functions and users in social or e-commerce networks exhibiting diverse interests. 
Tackling multi-label node classification (MLNC) on graphs has led to the development of various approaches. Some methods leverage graph neural networks (GNNs) to exploit label co-occurrence correlations, while others incorporate label embeddings to capture label proximity. However, these approaches fail to account for the intricate influences between labels in non-Euclidean graph data.
To address this issue, we decompose the message passing process in GNNs into two operations: \emph{propagation} and \emph{transformation}. 
We then conduct a comprehensive analysis and quantification of the influence correlations between labels in each operation. 
Building on these insights, we propose a novel model, Label Influence Propagation (\model). 
Specifically, we construct a label influence graph based on the integrated label correlations. 
Then, we propagate high-order influences through this graph, dynamically adjusting the learning process by amplifying labels with positive contributions and mitigating those with negative influence.
Finally, our framework is evaluated on comprehensive benchmark datasets, consistently outperforming SOTA methods across various settings, demonstrating its effectiveness on MLNC tasks\footnote{Our code is available at \url{https://github.com/Xtra-Computing/LIP_MLNC}.}.
\end{abstract}
% 